\documentclass[border=10pt]{standalone}
\usepackage{tikz}
\usepackage{circuitikz}
\ctikzset{logic ports=ieee}
\usetikzlibrary{positioning, calc, decorations.markings}

% Datapath of the if-then-else algorithm (if A = 5 then B := 8 else B := 13).
% Same building blocks and styles as lect_9_datapath_combine.

\tikzset{register/.style={muxdemux, muxdemux def={Lh=2.4, Rh=2.4, NL=2, NB=1, NT=1,w=4.8,NR=0},
                            muxdemux label={L1=Load, L2=Clock, R1=right1, R2=right,T1=$D_{3..0}$,B1=$Q_{3..0}$,},},
                 }

\begin{document}
\makeatletter\def\pgf@circ@font@tiny{\footnotesize}\makeatother
\begin{circuitikz}[
    font=\sffamily,
    label_text/.style={font=\small},
    bus/.style={-Latex, line width=1.5pt,
        postaction={decorate},
        decoration={markings, mark=at position 0.5 with {
            \draw[thick,-] (-2pt,-3pt) -- (2pt,3pt);
            \node[above left=0.5pt] {\small $4$};
        }}},
    bus plain/.style={-Latex, line width=1.5pt},
    control_signal/.style={-Latex, line width=0.5pt},
]

    % ---------------------------------------------------------------- registers
    \node[register] at (0,5) (A) {$A$};
    \node[register] at (6,5) (B) {$B$};

    % multiplexer for B: input 1 = '8' (left), input 0 = '13' (right)
    \draw[thick] (5.1,8.2) -- (6.9,8.2) -- (6.3,7.5) -- (5.7,7.5) -- cycle;
    \node[label_text] at (5.5,7.95) {$1$};
    \node[label_text] at (6.5,7.95) {$0$};
    \node[above,label_text] at (5.5,9.2) (c8) {$'8'$};
    \node[above,label_text] at (6.5,9.2) (c13) {$'13'$};
    \draw[bus plain] (c8) -- (5.5,8.2);
    \draw[bus plain] (c13) -- (6.5,8.2);
    \draw[bus] (6,7.5) -- (B.btpin 1);

    % input of A
    \node[above,label_text] at (0,9.2) {$Input$};
    \draw[bus] (0,9.2) -- (A.btpin 1);

    % ---------------------------------------------------------------- A = 5 detector
    % bus from A down to the detector, tapped by four inverted inputs
    \draw[bus] (A.bbpin 1) -- (0,1.8);
    \draw[line width=1.5pt] (0,1.8) -- (0,0.1);

    \node[and port, number inputs=4, rotate=180, anchor=out] (And) at (-2.7,1.0) {};
    % A = 5 = 0101: the inputs of the bits that are 0 (A3 and A1) are inverted
    % (rotated gate: in 1 = A0 ... in 4 = A3)
    \foreach \i in {1,2,3,4} {
        \draw (And.in \i) -- (0,0 |- And.in \i) node[circ, scale=0.6] {};
    }
    \foreach \i in {2,4} {
        \draw[fill=white] ($(And.in \i)+(-0.3,0)$) circle (2.6pt);
    }
    \node[label_text, anchor=south west] at (0.1,0 |- And.in 4) {$A_3$};
    \node[label_text, anchor=north west] at (0.1,0 |- And.in 1) {$A_0$};
    \draw[control_signal] (And.out) -- ++(-1.0,0) node[left, label_text] {$(A=5)$};

    % ---------------------------------------------------------------- tri-state output of B
    \node[buffer port, rotate=-90, yscale=-1, scale=0.5] (Tri) at (6,1.0) {};
    \draw[bus] (B.bbpin 1) -- (Tri.bin);
    \draw[bus plain] (Tri.bout) -- ++(0,-1.2) node[below, label_text] {$Output$};

    % ---------------------------------------------------------------- control signals
    \draw[control_signal] (-4,0 |- A.blpin 1) node[left, label_text] {$ALoad$} -- (A.blpin 1);
    \draw[control_signal] (3.0,0 |- B.blpin 1) node[left, label_text] {$BLoad$} -- (B.blpin 1);
    \draw[control_signal] (3.0,7.65) node[left, label_text] {$Muxsel$} -- (5.571,7.65);
    \draw[control_signal] (3.0,0 |- Tri.up) node[left, label_text] {$Out$} -- (Tri.up);

    % clock rail
    \draw[control_signal] (-4,3.5) node[left, label_text] {$Clock$} -- (4.2,3.5) |- (B.blpin 2);
    \draw[control_signal] (-2.7,3.5) to[short,*-] (-2.7,0 |- A.blpin 2) -- (A.blpin 2);

\end{circuitikz}
\end{document}
